\apendice{Anexo de sostenibilización curricular}

\section{Introducción a la Sostenibilidad}
Este anexo presenta una reflexión detallada sobre cómo el sistema AgroSkopos responde a los desafíos actuales de la sostenibilidad. Siguiendo las directrices de la Conferencia de Rectores de las Universidades Españolas (CRUE), el Trabajo de Fin de Grado no solo debe valorarse técnicamente, sino evaluando su impacto real. La transición hacia un modelo agrícola automatizado en entornos corporativos (B2B) ofrece una oportunidad inmejorable para integrar la ética profesional y la responsabilidad social corporativa. A continuación, se desgranan las tres dimensiones fundamentales de la sostenibilidad aplicadas a esta plataforma.

\section{Dimensión Medioambiental}
El pilar medioambiental constituye el núcleo fundamental de AgroSkopos. Históricamente, la agricultura a gran escala ha dependido de la sobreexplotación de recursos naturales y del uso intensivo de maquinaria de combustión fósil.

En primer lugar, el sistema fomenta la \textbf{Agricultura de Precisión}. Al recolectar datos a través de sensores (humedad, pH y NPK) y representarlos mediante mapas de calor, se proporciona a la empresa agrícola la inteligencia necesaria para aplicar el riego y los agroquímicos de forma localizada. Esta precisión previene uno de los problemas ecológicos más graves del sector: la escorrentía superficial de químicos que acaba contaminando los acuíferos subterráneos.

En segundo lugar, el proyecto impulsa la \textbf{Reducción de la Huella de Carbono}. El diseño de la plataforma está orientado a operar con vehículos 100\% eléctricos que reemplazan a los altamente contaminantes tractores diésel. La integración de la telemetría demuestra que es viable mantener el rendimiento agrícola operando con cero emisiones directas en la parcela de trabajo.

\section{Dimensión Económica}
La adopción de AgroSkopos por parte de empresas del sector primario supone una optimización sin precedentes de sus recursos operativos.

La adopción de la agricultura de precisión garantiza un impacto ecológico positivo y se traduce de forma inmediata en un ahorro de costes significativo. La reducción en la adquisición y el despliegue indiscriminado de insumos químicos (pesticidas y fertilizantes sintéticos) y agua de riego, permite a las empresas agrícolas aumentar sus márgenes de beneficio minimizando sistemáticamente el desperdicio.

Adicionalmente, la capacidad de automatizar misiones reduce enormemente la dependencia del trabajo manual para tareas repetitivas de exploración. Este factor permite a las organizaciones reestructurar sus equipos, derivando a los empleados hacia labores de mayor valor añadido: tareas de gestión analítica, toma de decisiones estratégicas y supervisión a través de la plataforma web.

\section{Dimensión Social y Ética Profesional}
La introducción de flotas de robots autónomos conlleva profundas implicaciones sociales. Lejos de suponer una amenaza para el empleo rural, sistemas como AgroSkopos actúan como un poderoso catalizador para la \textbf{modernización estructural y la mejora de las condiciones laborales}.

El trabajo agrícola tradicional exige un esfuerzo físico extremo y exposición continuada a factores de riesgo (insolación severa y manipulación de compuestos químicos tóxicos). Al delegar estas tareas en el vehículo autónomo, el perfil del trabajador humano asciende a un rol de operador tecnológico. De esta forma, se dignifica el trabajo de campo y se reducen drásticamente las tasas de accidentes laborales y enfermedades crónicas derivadas de la toxicidad química.

Desde la ética profesional de la ingeniería de software, la arquitectura incorpora validación estricta, gestión de usuarios basada en roles (RBAC) y mecanismos de seguridad activa (paradas de emergencia remotas), garantizando que las operaciones cumplan los estándares exigentes de seguridad corporativa.

\section{Alineación con los Objetivos de Desarrollo Sostenible (ODS)}
El diseño conceptual del proyecto AgroSkopos se alinea de forma transparente con la Agenda 2030 de las Naciones Unidas, impactando en los siguientes Objetivos de Desarrollo Sostenible:

\begin{itemize}
    \item \textbf{ODS 2 (Hambre Cero):} El *software* mejora sustancialmente la resiliencia y el rendimiento de los cultivos mediante una monitorización continua, ayudando a producir mayor cantidad de alimentos utilizando menor cantidad de recursos.
    \item \textbf{ODS 6 (Agua Limpia y Saneamiento):} Al ajustar con precisión las dosis de fertilizantes estrictamente necesarias, se previene directamente la filtración de nitratos dañinos, protegiendo las reservas hídricas locales.
    \item \textbf{ODS 7 (Energía Asequible y no Contaminante):} El uso de esta plataforma con flotas de robots eléctricos impulsa al sector primario a abandonar progresivamente los combustibles fósiles.
    \item \textbf{ODS 15 (Vida de Ecosistemas Terrestres):} La mitigación en el uso indiscriminado de pesticidas en grandes extensiones previene activamente la degradación del suelo arable, protegiendo la biodiversidad local.
\end{itemize}
